% ========================================================================
% فصل ۱۱: فناوری‌های نوظهور
% ========================================================================
\section{فناوری‌های نوظهور}
گسترش شبکه اقتضایی خودرویی برای سرویس‌های سامانه‌های حمل‌ونقل هوشمند، بهبود تأخیر تحویل داده، عملکرد پایدار شبکه، مقیاس‌پذیری و انعطاف‌پذیری را ضروری می‌کند \cite{hozouriOverviewVANETVehicular2023}. با توجه به روش‌های هوشمند مطرح‌شده در فصل پیش، فناوری‌های نوظهوری که به این نیازها پاسخ می‌دهند، در این فصل در شش حوزه بررسی می‌شوند: رایانش ابری، لبه و مه؛ شبکه‌های نسل پنجم و ششم؛ برش شبکه و مجازی‌سازی؛ شبکه‌های خودرویی نرم‌افزارمحور؛ معماری‌های ترکیبی؛ و دوقلوهای دیجیتال و شبیه‌سازی بلادرنگ.

\subsection{رایانش ابری، لبه و مه}
نخستین حوزه، بازتوزیع پردازش و ذخیره‌سازی از خودرو به ساختارهای رایانش ابری\LTRfootnote{\lr{Cloud Computing}} و لبه‌ای است. در این چارچوب، هر خودرو گره‌ای قدرتمند است و خوشه‌ای از آن‌ها می‌تواند به‌عنوان مزرعه‌ای از رایانه‌های فرا‌کن متحرک عمل کند؛ به این شبکهٔ نوین، ابر خودرویی\LTRfootnote{\lr{VANET Cloud}} گفته می‌شود \cite{hozouriOverviewVANETVehicular2023}. محاسبات تعاونی میان خودروها و ابر خودرویی از راه‌حل‌های پردازشی مرور کلان‌دادهٔ اینترنت خودروهاست و پیش‌بینی جریان ترافیک و مسیر کوتاه‌مدت را نیز برای بهینه‌سازی انتقال داده در آن مطرح شده است \cite{xuInternetVehiclesBig2018}. پیش‌بینی جریان ترافیک با یادگیری عمیق در ابر خودرویی نیز از جمله کارهای بررسی‌شده در مرور چالش‌های امنیتی است \cite{farsimadanReviewSecurityChallenges2025}. با این حال، روش‌های دفاعی سنتی عمدتاً به‌دلیل ماهیت متمرکزشان، که در آن سرورهای ابری نقش گلوگاه دارند، ضعیف و ناموفق‌اند \cite{farsimadanReviewSecurityChallenges2025}. هوش مصنوعی در برابر حملات سایبری آسیب‌پذیر است و ترکیب آن با زنجیره بلوکی، اینترنت خودروها را غیرمتمرکز، هوشمند و امن می‌کند، اما بار محاسباتی زنجیره بلوکی را افزایش می‌دهد؛ از این‌رو به معماری‌ای نیاز است که هوش مصنوعی و زنجیره بلوکی را با رایانش لبه\LTRfootnote{\lr{Edge Computing}} خودرویی ترکیب کند \cite{hozouriOverviewVANETVehicular2023}. رایانش لبه‌ای \lr{VANET} نیز از جمله موضوعات پژوهشی آینده در این مرور فهرست شده است \cite{hozouriOverviewVANETVehicular2023}.

\subsection{شبکه‌های نسل پنجم و ششم}
هم‌زمان با تحول معماری‌های محاسباتی، فناوری‌های دسترسی نیز در حال ورود نسل‌های جدید به این شبکه‌ها هستند. پهنای باند \lr{DSRC} محدود است و نمی‌تواند ارتباطات آیندهٔ \lr{V2X} را پشتیبانی کند؛ \lr{LTE-V}، \lr{C-V2X} و \lr{5G-NR} جایگزین‌های مطرح‌شده‌اند \cite{farsimadanReviewSecurityChallenges2025}. پروتکل لایهٔ پیوند در ارتباطات \lr{V2X} همچنان \lr{802.11p} است و با افزایش تقاضا برای تأخیر بسیار کم و قابلیت اطمینان بالا، مدل‌های مدیریت امنیتی سنتی بدون تحمیل هزینه و سربار عملیاتی قابل‌توجه، نیازها را برآورده نمی‌کنند \cite{farsimadanReviewSecurityChallenges2025}. در معماری‌های امنیتی نسل پنجم و ششم، مؤلفه‌هایی مانند \lr{AUSF}، \lr{SCMF}، \lr{ARPF}، \lr{PCF}، \lr{SEAF} و \lr{5G PPP} معرفی شده‌اند \cite{farsimadanReviewSecurityChallenges2025}. \lr{V2X} مبتنی بر نسل ششم نیز از جمله موضوعات پژوهشی آینده در یکی از مرورهای این فصل است \cite{hozouriOverviewVANETVehicular2023}.

\subsection{برش شبکه و مجازی‌سازی}
در بستر این فناوری‌های دسترسی، مدیریت منابع شبکه از راه برش شبکه و مجازی‌سازی نیز مطرح است. برش شبکه\LTRfootnote{\lr{Network Slicing}} و ایزولاسیون، از راهکارهای دفاعی در ارتباط خودرو با زیرساخت (\lr{V2I}) بررسی شده‌اند \cite{farsimadanReviewSecurityChallenges2025}. در لایهٔ سلولی، چالش‌هایی مانند نبود حفاظت رمزنگاری سیگنالینگ \lr{LTE}، استراق سمع، پارازیت و ایستگاه‌های بی‌مجوز مطرح‌اند؛ در برابر آن‌ها، شبکهٔ نرم‌افزارمحور\LTRfootnote{\lr{Software-Defined Networking, SDN}} و مجازی‌سازی توابع شبکه\LTRfootnote{\lr{Network Functions Virtualization, NFV}} از راهکارهای پیشنهادی هستند \cite{farsimadanReviewSecurityChallenges2025}.

\subsection{شبکه‌های خودرویی نرم‌افزارمحور}
مجازی‌سازی، مسیر را به معماری‌های نرم‌افزارمحور باز می‌کند. در معماری شبکه‌های خودرویی نرم‌افزارمحور\LTRfootnote{\lr{Software-Defined Vehicular Network, SDVN}}، صفحهٔ کنترل\LTRfootnote{\lr{Control Plane}} از صفحهٔ داده\LTRfootnote{\lr{Data Plane}} تفکیک می‌شود \cite{hozouriOverviewVANETVehicular2023}. کنترل‌کننده داده‌های خودروها، مانند موقعیت، سرعت و وضعیت اتصال، را جمع‌آوری می‌کند و بر پایهٔ \lr{GPS} تصمیم مسیریابی می‌گیرد \cite{hozouriOverviewVANETVehicular2023}. در صفحهٔ داده، خودروها و \lr{RSU}ها دارای سوئیچ \lr{OpenFlow} هستند؛ خودروها اجزای بنیادین و \lr{RSU}ها اجزای ثابت صفحهٔ داده‌اند و \lr{RSU} جدول جریان \lr{OpenFlow} را نگه می‌دارد؛ \lr{OpenFlow} نیز شناخته‌شده‌ترین رابط جنوبی است \cite{hozouriOverviewVANETVehicular2023}. اینترنت خودروهای نرم‌افزارمحور\LTRfootnote{\lr{Software-Defined IoV}} نیز به‌عنوان گسترشی از این رویکرد در مرور کلان‌داده مطرح شده است \cite{xuInternetVehiclesBig2018}. یادگیری تقویتی\LTRfootnote{\lr{Reinforcement Learning}} نیز برای تعیین سیاست پیوند ارتباطی امن در برابر خودروهای بدخواه، در شبکه‌های خودرویی نرم‌افزارمحور به کار رفته است \cite{farsimadanReviewSecurityChallenges2025}.

\subsection{معماری‌های ترکیبی}
هم‌پوشانی این معماری‌ها، دیدگاهی ترکیبی را برای طراحی امنیتی مطرح می‌سازد. ارتباطات \lr{V2X} در عمل درگیر ذی‌نفعان چندگانه‌ای است: گردانندگان شبکه، مراجع جاده‌ای، تولیدکنندگان تجهیزات اصلی، شهرداری‌ها و تأمین‌کنندگان خدمات \cite{farsimadanReviewSecurityChallenges2025}. از فناوری‌های دفتر کل توزیع‌شده\LTRfootnote{\lr{Distributed Ledger Technology}} مانند زنجیره بلوکی می‌توان برای تعامل‌پذیری قابل‌اعتماد میان این مشارکت‌کنندگان چندگانه بهره گرفت؛ هرچند این امر می‌تواند پیچیدگی زمانی اضافه کند \cite{farsimadanReviewSecurityChallenges2025}. این ملاحظه با نگرانی‌های مقیاس‌پذیری دفترکل‌های توزیع‌شده که در فصل ۸ مطرح شد، هم‌راستاست \cite{farsimadanReviewSecurityChallenges2025}.

\subsection{دوقلوهای دیجیتال و شبیه‌سازی بلادرنگ}
در نهایت، دارایی‌ها و هویت‌های دیجیتالی، حوزهٔ بعدی این فناوری‌ها را شکل می‌دهند. در تنگل \lr{IOTA}، دارایی‌های دیجیتالی‌ای وجود دارند که توسط تنگل اصلی و بدون کارمزد امن‌سازی می‌شوند و امکان «دوقلوسازی دیجیتالی یا توکن‌سازی\LTRfootnote{\lr{Tokenization}} رایگان هر سامانه یا دارایی» را فراهم می‌کنند \cite{sealeyIOTATangle202022}. نشانه‌های توکن نایب‌تعویض‌ناپذیر\LTRfootnote{\lr{Non-Fungible Token, NFT}} نیز برای اثبات احراز اصالت و مالکیت دستگاه‌های \lr{IoT} به کار می‌روند \cite{sealeyIOTATangle202022}. دوقلوی دیجیتال\LTRfootnote{\lr{Digital Twin}} در منابع مورد بررسی فقط در همین بافت عمومی \lr{IoT} مطرح شده است؛ پیوند آن با شبکه‌های خودرویی و شبیه‌سازی بلادرنگ در هیچ‌یک از این منابع وارد نشده است.
